\section{Fall Cleaning}

\begin{quotex}
It is then, as appears, the greatest of all lessons to know one's self. For if one knows himself, he will know God; and knowing God, he will be made like God ... God alone is in need of nothing, and rejoices most when He sees us bright with the ornament of intelligence.
\flright{\textsc{St Clement of Alexandria}}
\end{quotex}


I decided to do some fall cleaning this week. Not much got clean, but I did raise a lot of dust, enough to plague me with sniffles and allergies all night. I found an old ambient dub CD; I remember I used to listen to it often, but now I have no idea why. The Moroccan folk music CD is OK, nice beat, but you can't belly dance to it. In a book I read at the same time, I found a lock of hair. For some reason, it seemed important to me, the only proof that I had previously existed in the flesh. I tried to salvage it as some sort of relic, but a puff of wind blew it away. I also found some old notes scattered across several notepads (I spent most of the time reading them), which I will now try to clear away.

\paragraph{Knowing Thyself}


There are three ways to know thyself:

\begin{enumerate}


\item \textbf{As a physical object}. You fall from a tree at a certain rate, your blood pressure is nnn, etc.

\item \textbf{As a psychological object}. You know things about yourself:


  \begin{enumerate}

  
\item What you cherish
  
\item What you hate
  
\item What you fear
  \end{enumerate}


\item \textbf{As pure spirit}. ~You know yourself as subject, not as object.

That makes you responsible for what you think and what you are. This is the most difficult to get across.
\end{enumerate}


\paragraph{Old Books}


There is a lot of emphasis on choosing the most influential books. Of course, there is a benefit to wrestling with the mind of the best and original thinkers. Yet I often find myself going back to derivative works. They are synthetic, not argumentative or analytical; they bring ideas together, unlike polemical works. That makes them very useful. It is like taking a leisurely stroll down a well-trodden path to the beach, rather than having to create a path yourself through the brush. Not that you shouldn't strike out on your own every now and then.

It is nice to find apples on the ground. Sometimes you need to climb up and pick your own. But best of all is to grow your own tree.

\paragraph{Ouija, the Demonic, and Children of God}


I listened to a discussion on a religious radio station inspired by the recent movie, \emph{Ouija}, about whether the Ouija board opens a path to the demonic. This assumes, unfortunately, that it is difficult to contact the demonic, when it is the demons who are more likely contacting, and influencing, you. How do you recognize them? By anything that inhibits self-knowledge, especially all those misconceptions.

On the same show, the moderator announced that we are all ``children of God.'' I don't know when that became a popular belief, but it is not some sort of birthright. First of all, he intended it as the idea that we are all passive children; this is repulsive to normal men.

\textbf{To be a son means to partake in the nature of the father. So to be a son of God is to be like Him: active, powerful, free, and creative.}

My position: It is common to fall under the influence of demons. It is rare to be Godlike.

\paragraph{Free Speech}


I found this passage marked in \textbf{Max Stirner}'s \emph{The Ego and his Own}. Published a century and a half ago, it demonstrates how little has changed:

\begin{quotex}
The clamor of the Liberals for freedom of the press runs counter to their own principle, their proper will. They will what they do not will; they wish, they would like. Hence it is too that they fall away so easily when once so-called freedom of the press appears; then they would like censorship. Quite naturally. The State is sacred even to them; likewise morals. They behave toward it only as ill-bred brats, as tricky children who seek to utilize the weaknesses of their parents. Papa State is to permit them to say many things that do not please him, but papa has the right, by a stern look, to blue-pencil their impertinent gabble. If they recognize in him their papa, they must in his presence put up with the censorship of speech, like every child.
\end{quotex}


The American Association of University Professors published this guideline on free speech\footnote{\url{http://www.aaup.org/report/freedom-expression-and-campus-speech-codes}}:

\begin{quotationx}
Freedom of thought and expression is essential to any institution of higher learning. Universities and colleges exist not only to transmit knowledge. Equally, they interpret, explore, and expand that knowledge by testing the old and proposing the new. This mission guides learning outside the classroom quite as much as in class, and often inspires vigorous debate on those social, economic, and political issues that arouse the strongest passions. In the process, views will be expressed that may seem to many wrong, distasteful, or offensive. Such is the nature of freedom to sift and winnow ideas.

On a campus that is free and open, no idea can be banned or forbidden. No viewpoint or message may be deemed so hateful or disturbing that it may not be expressed.
\end{quotationx}


Please don't believe it.

\paragraph{Absolute Idealism}


I decided to arrange my books by topic. Idealism used to be the default philosophy for educated men, at least from the time of Plato. Germany led a revival in the 19\textsuperscript{th} century; there were a several British philosophers, peaking with the Vedanta-like \textbf{Francis Bradley}; in Italy, there was \textbf{Giovanni Gentile} and then \textbf{Julius Evola}, although his philosophy is mostly neglected. Western-educated Indians see in philosophical idealism a line of thought that is most compatible with Eastern modes of thought. Of course, if you recall, \textbf{Rene Guenon} regarded Neoplatonism as part of a valid tradition.

I don't know about the other topics, but I think I could use a part time librarian.

\paragraph{Not So Bon Mots}


If you want to know what passing through the tollhouses feels like, where you are met with challenging tests that you cannot pass on your own, then establish an intimate relationship with a woman.

I am reading \emph{The Elegance of the Hedgehog}, based on a recommendation. I haven't finished it, but it seems to be about a female Cioran. There are not too many women who have that tragic sense of life.

One aim of Critical Theory is to recover the history of the classes that have had no voice. However, it is not so very hard to reconstruct the mentality of the shudra; it suffices to watch a few television comedies which I attempted recently. Apparently, the items of interest include masturbation, anal sex, getting stoned, and the results of \emph{Dancing with the Stars}. Then everyone gets along happily in the end.

Nature has changed its meaning. It used to be one's ideal to emulate, but now it means whatever occurs empirically. So whatever animals do is natural for humans to do. So now we imitate animals; in a bygone era, we were taught to emulate the saints or even to imitate Christ.

I watch sports because they always get the facts right and the commentators are very knowledgeable about the game. That is the complete opposite of news shows.

I watch some TV to get a feel for popular culture. While watching Fox news one night, my son happened to listen in, then asked why all the women on the show sounded so stupid. Maybe because he hadn't listened to the men. Fox shows no cleavage, but a lot of thigh, probably so viewers will be motivated to buy Viagra, one of their sponsors.

One commentator this weekend explained that, for certain voting blocks, facts and explanations are of no value, since it is ``perception'' alone that counts. That cries out for discussion, but it was passed over by the panel. Apparently, it is known and considered quite normal.

Naturalists now consider it unwise to interfere with natural processes. That is why we don't feed the wild animals and we let forest fires run their course. That is not applied to human populations. Teilhard de Chardin wonders:

\begin{quotex}
How should we judge the efforts we lavish in all kinds of hospitals on saving what is so often no more than one of life's rejects?

To what extent should not the development of the strong (to the extent that we can define this quality) take precedence over the preservations of the weak?
\end{quotex}


Plato thought popular music was a danger to the order of society; was he ever wrong about anything?

Idea density and grammatical complexity are measures of mental ability. Specificity is another indicator, or at least it is a component of idea density.

Writers always leave something out, for the reader to discern.

I learned to think from \textbf{Benedict Spinoza}; I learned to understand social issues scientifically from \textbf{August Comte}. I learned holiness from \textbf{Miguel de Unamuno}.


\flrightit{Posted on 2014-11-02 by Cologero}

\begin{center}* * *\end{center}

\begin{footnotesize}
\begin{sffamily}
\texttt{David on 2014-11-03 at 00:37 said:}

Greetings. You speak of Idealism and of Plato. What about the critics of Aristotle regarding Plato ? Wouldn't we be just better of using and naming St Thomas since he went ahead and ''bested'' both of them ? Thanks.

\hfill

\texttt{Cologero on 2014-11-03 at 05:38 said:}

Yes, David, I've mentioned that many times and don't feel the need to repeat everything. Also, for me ``Plato'' is the ``representative type'' of a larger tradition. I don't like the word ``bested'' in this context, since philosophy (of the type mentioned) is not a battle, but part of an ongoing conversation. There are statics and there are dynamics. Therefore, thinking needs to expose the logical consequences of previous thinkers and integrate them into a larger whole.


\hfill

\texttt{Scardanelli on 2014-11-04 at 13:57 said:}

``But some of them like rare monks, perhaps as much as Dancing with the Stars''

Well of course... Yoga and Buddhism are high on the list of things modern women find super cool and attractive. Yoga is exercise and Buddhism is relaxing. If he were a saint, imitating Christ and preaching truths not of this world, I wonder if she would be as fascinated.

This seems to be the trouble with walking this path in the world... There is not only the difficulty of understanding, of meditating, of changing one's consciousness, but also of learning how to relate to those around you who are not on the same path. I am beginning to understand the novelty of celibacy.

\hfill

\end{sffamily}
\end{footnotesize}

